\clearpage
\ncuFrontTitle{符號說明}
\noindent
\begin{tabular}{@{}p{0.22\linewidth}p{0.72\linewidth}@{}}
\(x\) & 觀測樣本的輸入特徵向量。\\
\(y,\hat y\) & 實際標籤與預測標籤；關注類別為正類。\\
\(t_i,n\) & 第 \(i\) 筆樣本的觀測時間與資料筆數。\\
\(p_i(x)\) & 第 \(i\) 個基礎分類器對正類的預測機率。\\
\(S,|S|\) & 靜態集成的成員集合及其模型數量；統計比較式中的 \(S\) 另表示種子數。\\
\(p_S(x)\) & 子集 \(S\) 的等權集成機率。\\
\(\tau\) & 分類門檻。\\
\(\mathrm{TP},\mathrm{TN}\) & 真陽性與真陰性樣本數。\\
\(\mathrm{FP},\mathrm{FN}\) & 假陽性與假陰性樣本數。\\
\(\mathrm{TPR},\mathrm{FPR}\) & 真陽性率與假陽性率。\\
\(\mathrm{AUC}\) & ROC 曲線下面積。\\
\(\mathrm{AP}\) & 平均精確率。\\
\(\mathrm{Precision}\) & 正類預測的精確率。\\
\(\mathrm{Recall}\) & 正類召回率。\\
\(\mathrm{F1}\) & Precision 與 Recall 的調和平均數。\\
\(B\) & 統計推論的 bootstrap 重抽次數。\\
\(\Delta\) & 兩種模型在指定指標上的差值。\\
\end{tabular}
